% ============================================================================
% Fichier  : partie6/P06_C24_divulgation.tex
% Titre    : La Divulgation Coordonnée
% Auteur   : MINKA MI NGUIDJOÏ Thierry Emmanuel
% Version  : 1.0 -- novembre 2026
% ============================================================================
% Sources : ISO/IEC 29147 (divulgation de vulnérabilités) ; ISO/IEC 30111
% (traitement des vulnérabilités) ; loi camerounaise 2010/012. Cas d'une
% faille découverte dans un produit tiers : fictif.
% ============================================================================

\chapter{La Divulgation Coordonnée}
\label{chap:divulgation}
\chaptermark{La divulgation coordonnée}

\epigraph{%
    « Trouver une faille crée un devoir~:\\
    celui de la faire corriger sans nuire. »%
}{Principe du cours}

\niveauChapitre{Intermédiaire}{Chapitre~\ref{chap:droit}~;
Chapitre~\ref{chap:rapport}.}

\begin{boxobjectifs}
\begin{itemize}
    \item Définir la divulgation coordonnée de vulnérabilités.
    \item Distinguer divulgation privée, coordonnée et complète.
    \item Décrire le processus des normes ISO/IEC 29147 et 30111.
    \item Situer la divulgation dans le cadre juridique camerounais.
    \item Rédiger un signalement responsable.
\end{itemize}
\end{boxobjectifs}

\bigskip

% ============================================================================
\section*{Question directrice : que faire d'une faille découverte dans un logiciel utilisé par beaucoup d'autres~?}
% ============================================================================

En testant LiZbiZ, M.~TAMO découvre que la faille de l'ancien portail (T-09)
vient d'un composant répandu, utilisé par de nombreuses organisations. La
corriger chez LiZbiZ ne protège pas les autres. Comment signaler la faille à
l'éditeur sans exposer tout le monde~? C'est l'objet de la divulgation
coordonnée.

\medskip
\textbf{Plan du chapitre.}
\begin{enumerate}
    \item Trois façons de divulguer (section~\ref{sec:c24-trois}).
    \item Le processus coordonné (section~\ref{sec:c24-processus}).
    \item Les normes 29147 et 30111 (section~\ref{sec:c24-normes}).
    \item Le cadre juridique (section~\ref{sec:c24-droit}).
    \item Rédiger un signalement (section~\ref{sec:c24-signalement}).
\end{enumerate}

\separateur

% ============================================================================
\section{Trois Façons de Divulguer}
\label{sec:c24-trois}
% ============================================================================

\begin{tableaucours}{L{3.4cm}L{8.6cm}}{%
    \tableauHeader{Façon} &
    \tableauHeader{Principe et limite}
}
Privée & On informe le seul éditeur~; sûr, mais sans pression si l'éditeur
n'agit pas \\
Coordonnée & On informe l'éditeur, puis on publie après correction ou
délai convenu \\
Complète et immédiate & On publie tout aussitôt~; expose les utilisateurs
avant correctif \\
\end{tableaucours}
\vspace{-0.8em}
\begin{center}{\small\itshape Tableau~24.1. Trois modes de divulgation}\end{center}

La divulgation \termecle{coordonnée} est le compromis responsable~: elle
donne à l'éditeur le temps de corriger, puis informe les utilisateurs pour
qu'ils se protègent.

% ============================================================================
\section{Le Processus Coordonné}
\label{sec:c24-processus}
% ============================================================================

\begin{boxprocedure}
\textbf{Divulguer de façon coordonnée}
\begin{enumerate}
    \item Trouver le contact de sécurité de l'éditeur.
    \item Signaler la faille par un canal sûr, avec les éléments de
          reproduction.
    \item Convenir d'un délai de correction raisonnable.
    \item Laisser l'éditeur corriger et, le cas échéant, obtenir un
          identifiant public de vulnérabilité.
    \item Publier une description après le correctif, sans mode d'emploi
          d'attaque inutile.
\end{enumerate}
\end{boxprocedure}

\begin{boxalert}
\textbf{Violation typique.} Publier la faille avec son code d'exploitation
avant tout correctif, « pour forcer l'éditeur ».

\smallskip
\textbf{Règle~:} la divulgation vise à protéger, pas à nuire ni à se mettre
en avant. On laisse un délai raisonnable, on ne publie pas d'arme prête à
l'emploi, on pense d'abord aux utilisateurs exposés.
\end{boxalert}

% ============================================================================
\section{Les Normes 29147 et 30111}
\label{sec:c24-normes}
% ============================================================================

\begin{tableaucours}{L{4.0cm}L{8.2cm}}{%
    \tableauHeader{Norme} &
    \tableauHeader{Objet}
}
ISO/IEC 29147 & Comment \emph{recevoir} et publier les signalements de
vulnérabilités (côté celui qui divulgue et celui qui reçoit) \\
ISO/IEC 30111 & Comment \emph{traiter} en interne une vulnérabilité signalée
jusqu'à sa correction \\
\end{tableaucours}
\vspace{-0.8em}
\begin{center}{\small\itshape Tableau~24.2. Deux normes complémentaires~\cite{iso29147,iso30111}}\end{center}

Les deux normes sont complémentaires~: 29147 organise l'échange avec le
monde extérieur, 30111 organise le travail interne de correction.

% ============================================================================
\section{Le Cadre Juridique}
\label{sec:c24-droit}
% ============================================================================

\begin{boxjuris}
La découverte d'une faille ne vaut pas autorisation d'y accéder. Au
Cameroun, la loi 2010/012 sur la cybersécurité réprime l'accès et le
maintien frauduleux dans un système, même sans dommage
(Chapitre~\ref{chap:droit}). Un chercheur qui découvre une faille hors d'un
mandat doit s'en tenir à ce qu'il a constaté sans intrusion supplémentaire,
et signaler.
\end{boxjuris}

% ============================================================================
\section{Rédiger un Signalement}
\label{sec:c24-signalement}
% ============================================================================

\begin{boxmetier}{Pentesteur et RSSI}
\textbf{Signaler avec mesure.} Un bon signalement est factuel, reproductible
et courtois~: il décrit la faille, son impact, les conditions de
reproduction, et propose un délai. Il n'exige rien, il coopère.
\end{boxmetier}

\separateur

\begin{boxresume}
\textbf{Ce qu'il faut retenir de ce chapitre~:}
\begin{itemize}
    \item \textbf{Trouver une faille} crée le devoir de la faire corriger
          sans nuire.
    \item \textbf{La divulgation coordonnée} donne le temps de corriger, puis
          informe.
    \item \textbf{29147} organise l'échange, \textbf{30111} le traitement
          interne.
    \item \textbf{Découvrir n'autorise pas à accéder}~: la loi 2010/012
          s'applique.
    \item \textbf{Un signalement} est factuel, reproductible, coopératif.
\end{itemize}
\end{boxresume}

\begin{boxinfo}
\textbf{Points de contrôle de l'élève.} Je sais~: distinguer les trois modes
de divulgation~; dérouler le processus coordonné~; distinguer 29147 et
30111~; rappeler le cadre juridique~; rédiger un signalement.
\end{boxinfo}

\bigskip

\begin{boxexo}[title={\ding{45}\quad Exercice~24.1. Trois modes \niveaubase}]
Distinguez divulgation privée, coordonnée et complète immédiate.
\end{boxexo}

\begin{boxexo}[title={\ding{45}\quad Exercice~24.2. Normes \niveaubase}]
Quelle norme organise la réception des signalements~? Laquelle organise leur
traitement interne~?
\end{boxexo}

\begin{boxexo}[title={\ding{45}\quad Exercice~24.3. Devoir \niveaubase}]
Pourquoi corriger T-09 chez LiZbiZ ne suffit-il pas~?
\end{boxexo}

\begin{boxexo}[title={\ding{45}\quad Exercice~24.4. Délai \niveaumoyen}]
Pourquoi laisse-t-on un délai à l'éditeur avant publication~? Qu'est-ce qu'un
délai raisonnable~?
\end{boxexo}

\begin{boxexo}[title={\ding{45}\quad Exercice~24.5. Pas d'arme \niveaumoyen}]
Pourquoi ne publie-t-on pas le code d'exploitation prêt à l'emploi~?
\end{boxexo}

\begin{boxexo}[title={\ding{45}\quad Exercice~24.6. Droit \niveaumoyen}]
Un chercheur découvre une faille sur un site public. Que peut-il faire, que
ne peut-il pas faire, selon la loi 2010/012~?
\end{boxexo}

\begin{boxexo}[title={\ding{45}\quad Exercice~24.7. Contact \niveaumoyen}]
Comment trouve-t-on le contact de sécurité d'un éditeur~?
\end{boxexo}

\begin{boxexo}[title={\ding{45}\quad Exercice~24.8. Signalement \niveaumoyen}]
Rédigez le signalement de la faille de composant à l'éditeur.
\end{boxexo}

\begin{boxexo}[title={\ding{45}\quad Exercice~24.9. Politique \niveauexpert}]
Rédigez une politique de divulgation que LiZbiZ pourrait publier pour
recevoir des signalements, inspirée de 29147.
\end{boxexo}

\begin{boxexo}[title={\ding{45}\quad Exercice~24.10. Processus interne \niveauexpert}]
Décrivez le processus interne de traitement d'une vulnérabilité signalée,
inspiré de 30111.
\end{boxexo}

\bigskip

\begin{boxinfo}
\textbf{Pour aller plus loin.}
\begin{itemize}
    \item ISO/IEC 29147~\cite{iso29147}~; ISO/IEC 30111~\cite{iso30111}.
    \item Loi 2010/012~\cite{loi2010012}.
\end{itemize}
\end{boxinfo}

\begin{boxinfo}
\textbf{Transition.} Toutes ces compétences se cultivent dans un cadre
d'entraînement~: le Chapitre~\ref{chap:labo} traite du laboratoire et de la
pratique continue.
\end{boxinfo}
